% !PRINTED 599
\weberpage{599}
%<<CONT
Grad $h_r$ des Durchschnittes von $\Phi'$ und $\Psi_r$. Der Grad von $\Phi_r$
ist also gleich $f g n' : h_r$.

Wenn wir jetzt mit $\varphi_r$ alle Elemente von $\Phi_r$ bezeichnen, so
führen alle Primideale $\mathfrak{p}\,|\,\varphi_r$ in $\Omega$ zu demselben Primideale
$\mathfrak{p}_r'$ in $\Omega'$. Es ist aber noch zu zeigen, dass zwei verschiedene
dieser Systeme $\Phi_1, \Phi_2$ zu \gesperrt{verschiedenen} Primidealen $\mathfrak{p}_1', \mathfrak{p}_2'$
führen.

Zunächst ist ersichtlich, dass die sämmtlichen $\mathfrak{p}\,|\,\varphi_2$ von den
$\mathfrak{p}\,|\,\varphi_1$ verschieden sind. Denn wenn $\mathfrak{p}\,|\,\varphi_1 = \mathfrak{p}\,|\,\varphi_2$ ist, so ist
auch $\mathfrak{p} = \mathfrak{p}\,|\,\varphi_2\,\varphi_1^{-1}$, also $\varphi_2\,\varphi_1^{-1}$ in $\boldsymbol{\Psi}$ und folglich $\varphi_2$ in
$\boldsymbol{\Psi}\varphi_1$, d.\,h. in $\Phi_1$ enthalten, gegen die Voraussetzung. Wir
können also eine ganze Zahl $\alpha$ in $\Omega$ annehmen, die durch
$\mathfrak{p}\,|\,\varphi_1$, aber durch keines der Ideale $\mathfrak{p}\,|\,\varphi_2$ theilbar ist. Dann
ist auch keine der Zahlen $\alpha\,|\,\varphi'$ durch $\mathfrak{p}\,|\,\varphi_2$ theilbar, weil sonst
$\mathfrak{p}\,|\,\varphi_1\,\varphi' = \mathfrak{p}\,|\,\varphi_2$, also gegen die Voraussetzung $\varphi_2$ in $\Phi_1$
enthalten wäre. Das Product $\alpha'$ aller von einander verschiedener
$\alpha\,|\,\varphi'$, welches eine in $\Omega'$ enthaltene Zahl ist, ist zwar durch
$\mathfrak{p}_1'$, nicht aber durch $\mathfrak{p}_2'$ theilbar. Folglich ist $\mathfrak{p}_1'$ von $\mathfrak{p}_2'$
verschieden, und es ergiebt sich also, dass $e'$ und
nicht mehr \gesperrt{verschiedene} Primideale $\mathfrak{p}'$ in $\mathfrak{P}$ aufgehen.
Wir setzen daher

\weq{3}{\mathfrak{P} = \mathfrak{p}_1'^{\,a_1}\,\mathfrak{p}_2'^{\,a_2}\cdots\mathfrak{p}_{e'}'^{\,a_{e'}},}

worin die Exponenten $a_1, a_2, \ldots, a_{e'}$ noch näher zu bestimmende
natürliche Zahlen sind.

Um nun die Primideale $\mathfrak{p}_r'$ im Körper $\Omega$ zu zerlegen, können
wir die Resultate der \S.\,160, 162 benutzen, indem wir einfach
$\Omega'$ an Stelle von $R$ treten lassen. Bedeutet dann $X_r'$ den Durchschnitt
 von $\Phi'$ mit $X_r = \varphi_r^{-1}\,X\,\varphi_r$, so ist $X_r$ der Inbegriff aller % [sic] source omits the prime; the next sentence correctly refers to $X_r'$
% [sic: The source says "Körpers" although $\Phi'$ denotes the group, not the field.]
Substitutionen des Körpers $\Phi'$, die der Bedingung

\[
\tau\,|\,\chi_r' \equiv \tau \ (\mathrm{mod}\ \mathfrak{p}_r)
\]

genügen. Bezeichnen wir daher den Grad von $X_r'$ mit $g_r$, so ist
$\mathfrak{p}_r'$ durch $\mathfrak{p}_r^{g_r}$, aber durch keine höhere Potenz von $\mathfrak{p}_r$ theilbar.

Der Durchschnitt $\Psi_r'$ von $\Psi_r$ mit $\Phi'$, dessen Grad wir schon
oben mit $h_r$ bezeichnet haben, ist der Inbegriff aller Substitutionen
 $\psi_r'$ in $\Phi'$, die der Bedingung

\[
\mathfrak{p}_r\,|\,\psi_r' = \mathfrak{p}_r
\]

genügen, und wenn wir daher $\Phi'$ in die Nebengruppen

\weq{4}{\Phi' = \Psi_r'\varphi_{r,1}' + \Psi_r'\varphi_{r,2}' + \cdots + \Psi_r'\varphi_{r,e_r}} % [sic] source omits the prime on the final $\varphi$, unlike the two preceding terms

zerlegen, so ist

\weq{5}{n' = e_r\,h_r.}
